\setchapterpreamble[u]{\margintoc}
\chapter{Future Work}

One of the main improvements that could further strengthen this thesis would be the incorporation of a more detailed analysis aimed at distinguishing healthy from malignant cells. During the development of this project, the laboratory was still awaiting pathological annotations identifying whether individual cells were malignant or non-malignant. For this reason, such analyses were not integrated into the present work. Nevertheless, including this information would provide a more precise biological framework and could lead to more robust conclusions regarding tumour organisation and treatment response.

Importantly, the absence of visually detectable tumour tissue does not necessarily imply the complete absence of malignant transcriptional activity. Certain samples may contain residual disease characterised by epithelial populations that appear less morphologically malignant while still retaining transcriptional signatures associated with tumoural states. Investigating these transcriptional alterations in greater depth could provide valuable insights into tumour persistence, treatment resistance, and intermediate cellular states following therapy.

An important direction for future work would be the validation of the spatial patterns identified in this study in larger and independent MIBC cohorts treated with immunotherapy. While the present analyses revealed significant differences in the spatial distribution of cell populations between responders and non-responders, further studies are required to determine whether these patterns are consistently associated with treatment outcome. In particular, quantitative analyses of tumour--immune cell distances and immune-cell infiltration into tumour regions could help confirm whether the fibroblast-rich epithelial communities observed in non-responders correspond to a true immune-excluded phenotype. Additionally, a more detailed characterization of the fibroblast populations involved, including the identification of cancer-associated fibroblast subtypes and extracellular matrix-related programs according to their phenotype, would provide further insight into the mechanisms underlying treatment resistance.

Finally, the potential of ST data extends far beyond the analyses performed in this thesis. Future investigations could integrate spatial transcriptomics with additional molecular and clinical biomarkers to improve patient stratification prior to treatment selection. Analyses of ligand--receptor interactions, immune-related signalling pathways, transcription factor activity, and other tumour microenvironment features could reveal additional mechanisms associated with treatment response. Together with the spatial biomarkers identified in this work, these approaches may contribute to the development of more accurate predictive models capable of identifying which patients are most likely to benefit from immunotherapy and which may require alternative therapeutic strategies, thereby advancing personalized medicine in MIBC.


%Higher Cohen’s D (dark green color) indicates a higher relative abundance of a given cell type in each community in responders, whereas a higher value towards the orange color indicates the opposite